\section{Extending ViZDoom's Research Scope}
\label{sec:updates}
The additions since the historical platform broaden both the evidence available to an agent and the structure of the tasks it can attempt. Table~\ref{tab:updates} connects each addition to a research question. We distinguish new access to information from inherited engine capabilities: object labels, full-map loading, and multiplayer predate this update, while semantic categories and the newer evaluation interfaces extend their usefulness.

\begin{table*}[t]
\caption{Extensions to \vzd{} and the research questions they enable. The MCP server is the interface for the proposed language-model benchmark; the PettingZoo wrapper is a separate multi-agent integration under development.}
\label{tab:updates}
\small
\begin{tabularx}{\textwidth}{@{}P{0.17\textwidth}P{0.34\textwidth}Y@{}}
\toprule
Extension & Added capability & Research question \\
\midrule
Audio buffer & Access to game sounds alongside the visual stream & When does hearing improve decisions under partial observability? \\
Text buffer & Machine-readable feedback about in-game events & Can agents ground event messages in their actions and remembered state? \\
Label categories & Semantic roles associated with visible objects & Can semantic supervision improve perception and transfer? \\
Full-level environments & A common route to complete Doom-level objectives & Can agents compose local skills into sustained, sparse-reward behavior? \\
PettingZoo wrapper & Standardized interaction with multiple agents & How well do learned strategies transfer across opponents and roles? \\
MCP server & Observation, action, and selective visual-memory tools & How do model capability and harness design affect closed-loop performance? \\
\bottomrule
\end{tabularx}
\end{table*}

\subsection{Audio and Text for Action}
The audio buffer supplements the visual observation with sounds generated by the game. An agent may hear an event that occurs outside its current view or use a sound to distinguish states that look similar. This makes audiovisual control a natural way to study partial observability: the relevant question is whether an additional modality changes a decision, rather than whether a model can recognize the sound in isolation. The standard Basic Audio task provides a controlled example involving an invisible, sound-emitting target.

The text buffer supplies recent in-game notifications. These messages can communicate the outcome of an interaction or a scripted event without requiring the agent to read text from an image. Basic Notifications, for example, uses text to identify which target should be attacked. Such tasks connect language understanding to consequential action while retaining a visually grounded control problem.

Textual feedback is event information, not a general narration of the world. Its usefulness depends on what the scenario communicates, whether that information is already visible, and whether the agent retains it until needed. Comparisons should therefore distinguish the effect of additional information from the effect of presenting the same information in a more accessible form. Audio removal, delayed messages, and matched visual cues provide interventions for studying these differences.

\subsection{Semantic Representations}
Semantic categories extend existing object labels with broader roles, such as enemies, weapons, and players. They support questions about whether agents can learn representations that transfer across changes in appearance while preserving an object's functional role. Category supervision can also connect visual learning to the distinctions that matter for control: avoiding a threat, collecting a resource, or identifying another player.

Two uses deserve separate evaluation. Categories may supervise a perception model during training, with only images supplied at test time; alternatively, category information may be given directly to the acting agent. The latter measures performance with privileged perception and changes the task. Categories are therefore useful both as learning targets and as a diagnostic intervention, but gains in these settings have different interpretations. Custom objects may also require checking that the assigned categories match the intended semantics.

\subsection{Composing Skills in Complete Levels}
Full-level environments provide a common interface to objectives extending beyond the compact scenarios traditionally used in \vzd{}. Doom and Doom~II levels combine navigation and combat with exploration, resource management, and interactions with doors, switches, and keys. The addition is a more accessible route to these tasks; compatibility with Doom maps itself is an inherited capability. Independently reset levels also remain distinct from campaigns that carry inventory and other state across maps.

A terminal completion reward exposes the difficulty of discovering and executing useful intermediate behaviors without dense external guidance. A sequence such as finding a key, remembering a locked door, returning to it, and reaching the exit illustrates the interaction between exploration, memory, and delayed credit. Not every level requires this sequence. The opening maps offer recognizable reference tasks, while stronger claims about long-horizon composition require levels with verified dependencies and backtracking.

Full-level evaluation can thus connect success on isolated scenarios to success on composed objectives. A model that aims reliably but fails to complete a level may have difficulty deciding where to go, managing resources, or recovering from an incorrect plan. These distinctions are more informative than interpreting every failure as a deficit in low-level control.

\subsection{Interaction with Other Agents}
The PettingZoo wrapper connects native multiplayer Doom to a standard multi-agent interface. It supports research on self-play, opponent adaptation, and role-dependent behavior in a shared first-person world. The integration is being developed separately from the main library; its current observation interface is visual, so multimodal support in the core should not be assumed to extend automatically to multiplayer evaluation.

Multi-agent tasks add a source of uncertainty that single-player levels do not capture: other policies change their behavior in response to the agent. Evaluating against held-out opponents and exchanging player roles can reveal whether a strategy is robust or depends on a narrow training interaction. This complements the full-level study, where uncertainty primarily comes from incomplete information about the environment.
